\subsection{Purpose}
The purpose of Best Bike Paths (BBP) is to provide cyclists with a dedicated platform to discover, share and evaluate cycling routes, while allowing them to manage and track their own trips. The system aims to improve the safety, comfort, and overall experience of cycling by offering information on route conditions, obstacles, and environmental factors. Additionally, BBP create a community of cyclists who can exchange knowledge, follow high-quality routes and monitor personal progress over time.
\subsection{Scope}
The design of Best Bike Paths (BBP) focuses on defining the system’s structure, components, and interactions necessary to implement its functionalities. This document covers the design of the main modules that enable cyclists to record, manage and share bike routes.
The system includes user management features, allowing registration, authentication, and profile modification for those who choose to create an account. It allows route creation through both manual creation and automatic tracking via GPS and mobile device sensors, storing relevant statistics such as distance, time, and average speed. Routes are enriched with additional information, including environmental and meteorological data collected from external services.\\
BBP also incorporates a social component, enabling users to publish routes, explore public paths, and contribute to community ratings. Additionally, the system provides search and navigation functionality, suggesting paths based on quality scores, estimated travel time and overall status.\\
The design covers the system’s four-tier architecture, detailing the presentation, routing tier, application logic and data layers, as well as the interactions between them that ensure modularity, scalability and distribution. 
\subsection{Definitions, Acronyms, Abbreviations}
\subsubsection{Definitions}
\begin{itemize}
    \item \textbf{Guest:} The term refers to the unregistered user. More details are provided later in the document.
    \item \textbf{Cyclist:} The term refers to the registered user on the platform. More details are provided later in the document.
    \item \textbf{Path:} A path is a set of streets rideable by the cyclist from a starting point to an end point. It is automatically created when a cyclist records a route for the first time on a path that does not exist. Once a path is created, it behaves as a route's family aggregating all similar routes under the same path.
    \item \textbf{Route:} A route can be manually created by the cyclist or automatically recorded by the system. Each route is related to a specific path and contains additional details that describe the actual performance of the cyclist on that specific path.
    \item \textbf{Path's Statistics:} Metrics that describe a path: average time, average distance, status and score.
    \item \textbf{Route's Statistics:} Metrics that describe a route: date, duration, distance, average speed, weather conditions and status.
    \item \textbf{Score:} A metric that measures the quality of a path. It is calculated by the system and is based on the status of the path and how effective the path is in terms of time. 
    \item \textbf{Status:} A metric that reflects the quality and condition of the road (e.g., optimal, medium, sufficient, requires maintenance, bad).
    \item \textbf{Tolerance:} It is a parameter for the search that specifies how far the suggested routes can deviate from the selected points, allowing greater flexibility.
    \item \textbf{Tier:} It is a layer inside the architecture of the system.
\end{itemize}
\subsubsection{Acronyms}
\begin{itemize}
    \item \textbf{BBP:} Best Bike Paths
    \item \textbf{RASD:} Requirement Analysis and Specification Document
    \item \textbf{DD:} Design Document
    \item \textbf{UI:} User Interface
    \item \textbf{API:} Application Programming Interface
    \item \textbf{DBMS:} Database Management System
    \item \textbf{DTO:} Data Transfer Object
\end{itemize}
\subsubsection{Abbreviations}
\begin{itemize}
    \item \textbf{I:} Interface
    \item \textbf{R:} Requirement
\end{itemize}
\subsection{Revision History}
08/01/2026 - Updates on data types on section 2.5 Component interface and minor updates to the sequence diagrams.
\subsection{Reference Documents}
The assignment for this document and all the information included refer to the following documentation:
\begin{itemize}
    \item The specification for the 2025/26 Requirement Engineering and Design Project for the Software Engineering II course.
    \item The slides on the WEBEEP page of the Software Engineering II course.
\end{itemize}
\subsection{Document Structure}
\begin{enumerate}
    \item \textbf{Introduction}: provides a concise overview of the project, outlining the motivations behind its development and the goals to be achieved, as already introduced in the RASD;
    \item \textbf{Architectural Design}: describes the system architecture at different levels of abstraction, highlighting its structure and the relationships among its components;
    \item \textbf{User Interface Design}: shows the structure and layout of the platform’s user interface;
    \item \textbf{Requirements Traceability}: explains how the goals identified in the RASD are fulfilled through the mapping between requirements and the components described in this document;
    \item \textbf{Implementation, Integration and Test Plan}: details the strategies and activities related to implementation, integration, and testing, providing guidance for developers during the development process;
    \item \textbf{Effort Spent}: reports the amount of time each group member dedicated to the preparation of this document;
    \item \textbf{References}: lists the documents, standards, and tools consulted for the drafting of this document. 
\end{enumerate}